\documentclass[11pt]{article}
\usepackage{fontspec}
\setmainfont[
  Path=fonts/,
  Extension=.ttf,
  UprightFont=*-400,
  BoldFont=*-700
]{NotoSerifSC}
\setsansfont[
  Path=fonts/,
  Extension=.ttf,
  UprightFont=*-400,
  BoldFont=*-700
]{NotoSerifSC}
\XeTeXlinebreaklocale "zh"
\XeTeXlinebreakskip=0pt plus 1pt
\usepackage[a4paper,margin=1.70cm]{geometry}
\usepackage{amsmath,amssymb,mathtools,bm}
\usepackage{booktabs,longtable,array}
\usepackage{enumitem}
\setlength{\parindent}{2em}
\setlength{\parskip}{0.35em}
\setlist{nosep,leftmargin=2em}
\allowdisplaybreaks[3]
\numberwithin{equation}{section}
\pagestyle{plain}

\newcommand{\R}{\mathbb{R}}
\newcommand{\E}{\mathbb{E}}
\newcommand{\Pp}{\mathbb{P}}
\newcommand{\1}{\mathbf{1}}
\newcommand{\T}{\mathsf{T}}
\newcommand{\F}{\mathrm{F}}
\newcommand{\cN}{\mathcal{N}}
\newcommand{\cL}{\mathcal{L}}
\newcommand{\cS}{\mathcal{S}}
\newcommand{\cH}{\mathcal{H}}
\newcommand{\cA}{\mathcal{A}}
\newcommand{\cM}{\mathcal{M}}
\newcommand{\cT}{\mathcal{T}}
\newcommand{\diag}{\operatorname{diag}}
\newcommand{\tr}{\operatorname{tr}}
\newcommand{\Var}{\operatorname{Var}}
\newcommand{\Cov}{\operatorname{Cov}}
\newcommand{\softmax}{\operatorname{softmax}}
\newcommand{\KL}{D_{\mathrm{KL}}}
\newcommand{\sg}{\operatorname{sg}}
\begin{document}
    \begin{center}
      {\LARGE\bfseries 39\_Continuous Thought Machines}
    \end{center}
    \vspace{0.8em}

    \section{符号与维度}
    \begin{longtable}{>{$}l<{$} >{$}l<{$} p{10cm}}
    \toprule
    \text{符号} & \text{维度} & \text{含义}\\
    \midrule
    a_t,z_t & \R^D & 内部时刻的预激活与神经元状态。\\
        A_t & \R^{D\times L} & 每个神经元的时间历史。\\
        Z_t & \R^{D\times L_s} & 同步窗口状态矩阵。\\
        S_t & \R^{D\times D} & 神经同步矩阵。\\
        r_{ij} & \R_{\ge0} & 神经元对的时间衰减率。\\
    \bottomrule
    \end{longtable}

    所有向量默认是列向量；未特别说明时，期望同时对数据、噪声和采样时刻取值。下文只展开论文中承担方法逻辑的公式，并把所需假设写在推导发生的位置。

\section{共享递归块的状态与梯度}

设同一参数化块被递归调用 $T$ 次：
\begin{equation}
 h_{t+1}=F_\theta(h_t,x),\qquad t=0,\ldots,T-1.
\end{equation}
令末端损失 $\cL(h_T)$。对状态的反向递推
\begin{equation}
 g_t=\frac{\partial\cL}{\partial h_t}
 =J_{h_t}F_\theta(h_t,x)^Tg_{t+1}.
\end{equation}
因此
\begin{equation}
 g_t=\left(\prod_{j=t}^{T-1}J_{h_j}F_\theta^T\right)g_T,
\end{equation}
其中乘积按时间反序作用。共享参数在每次调用都贡献梯度：
\begin{equation}
 \boxed{\nabla_\theta\cL
 =\sum_{t=0}^{T-1}
 J_\theta F_\theta(h_t,x)^Tg_{t+1}.}
\end{equation}
参数量不随 $T$ 增长，但激活存储和计算量通常线性增长；共享不等于免费。

若 $\|J_{h_t}F\|_2\le\rho$，则
\begin{equation}
 \|g_t\|_2\le\rho^{T-t}\|g_T\|_2.
\end{equation}
$\rho<1$ 导致长程梯度衰减，$\rho>1$ 可能爆炸。残差形式
$F(h)=h+G(h)$ 的 Jacobian 为 $I+J_G$，给出直接通路但仍需控制谱范数。

\section{动态深度的期望计算量}

令 token $i$ 的递归深度为离散变量
$D_i\in\{1,\ldots,T\}$，概率 $p_{it}=\Pp(D_i=t)$。期望深度
\begin{equation}
 \E[D_i]=\sum_{t=1}^{T}tp_{it}.
\end{equation}
也可用生存概率写成
\begin{equation}
 \E[D_i]=\sum_{t=1}^{T}\Pp(D_i\ge t).
\end{equation}
证明为
\begin{align}
 \sum_{t=1}^{T}\Pp(D_i\ge t)
 &=\sum_{t=1}^{T}\sum_{d=t}^{T}p_{id}\\
 &=\sum_{d=1}^{T}\sum_{t=1}^{d}p_{id}
 =\sum_ddp_{id}.
\end{align}
若每次递归单 token 代价为 $c_i$，批次期望计算
\begin{equation}
 \E[C]=\sum_i c_i\E[D_i].
\end{equation}
预算正则可写
\begin{equation}
 \cL_{\rm budget}=\lambda
 \left(\frac1n\sum_i\E[D_i]-D_{\rm target}\right)^2.
\end{equation}
其对路由 logit 的梯度会把所有 token 耦合，因为批平均出现在同一平方项中。

\section{离散路由的梯度估计}

若 $D\sim p_\phi(D\mid x)$，期望损失
\begin{equation}
 J(\phi)=\sum_Dp_\phi(D\mid x)\cL_D
\end{equation}
的 score-function 梯度
\begin{align}
 \nabla_\phi J
 &=\sum_Dp_\phi(D)\cL_D\nabla_\phi\log p_\phi(D)\\
 &=\E[(\cL_D-b(x))\nabla_\phi\log p_\phi(D)],
\end{align}
其中基线不改变期望，因为
$\E\nabla\log p=\nabla\sum_Dp(D)=0$。该估计无偏但方差可能高。

Gumbel--Softmax 松弛令
\begin{equation}
 y_t=\frac{\exp((\ell_t+g_t)/\tau)}
 {\sum_j\exp((\ell_j+g_j)/\tau)},\qquad
 g_t=-\log[-\log u_t],\quad u_t\sim U(0,1).
\end{equation}
$\tau\to0$ 时趋向 one-hot，但 Jacobian 含 $1/\tau$，梯度可能尖锐。
直通版本前向取 $\operatorname{onehot}(\arg\max y)$、反向用软 $y$ 的梯度，因此是有偏估计。

\section{同步统计量的矩阵表达}

设神经元内部时间序列 $z_i(1),\ldots,z_i(T)$，去均值向量
\begin{equation}
 \widetilde z_i=z_i-\bar z_i\1,\qquad
 \bar z_i=\frac1T\sum_tz_i(t).
\end{equation}
归一化相关
\begin{equation}
 S_{ij}=\frac{\widetilde z_i^T\widetilde z_j}
 {\sqrt{\|\widetilde z_i\|_2^2+\epsilon}
  \sqrt{\|\widetilde z_j\|_2^2+\epsilon}}
\end{equation}
形成同步矩阵 $S$。若把标准化后的行堆成 $Z\in\R^{m\times T}$，则
\begin{equation}
 S=ZZ^T.
\end{equation}
所以 $S$ 半正定、$\operatorname{rank}(S)\le T$。对任意 $a$，
\begin{equation}
 a^TSa=\|Z^Ta\|_2^2\ge0.
\end{equation}
这类二阶表征包含神经元两两时间关系，计算和存储却为 $O(m^2T)$ 与 $O(m^2)$。

\section{指数衰减同步的在线递推}

对乘积统计定义
\begin{equation}
 C_t=\lambda C_{t-1}+(1-\lambda)z_tz_t^T,
 \qquad0\le\lambda<1.
\end{equation}
递归展开
\begin{equation}
 C_t=\lambda^tC_0+(1-\lambda)
 \sum_{k=1}^{t}\lambda^{t-k}z_kz_k^T.
\end{equation}
历史样本相隔 $h$ 步的权重为 $(1-\lambda)\lambda^h$，半衰期
\begin{equation}
 h_{1/2}=\frac{\log(1/2)}{\log\lambda}.
\end{equation}
若 $\lambda$ 接近 $1$，$\log\lambda\approx-(1-\lambda)$，
故 $h_{1/2}\approx\log2/(1-\lambda)$。

若对每一步都保留递归图，梯度穿过
\begin{equation}
 \frac{\partial C_t}{\partial C_{t-k}}=\lambda^kI,
\end{equation}
长程梯度指数衰减。若在线统计使用停止梯度，则它作为状态提供历史信息，却不承担跨全部时间的反向传播。

\section{固定点与隐式微分}

当递归收敛到固定点 $h^*=F_\theta(h^*,x)$，对参数微分：
\begin{equation}
 dh^*=J_hF\,dh^*+J_\theta F\,d\theta.
\end{equation}
若 $I-J_hF$ 可逆，
\begin{equation}
 \frac{\partial h^*}{\partial\theta}
 =(I-J_hF)^{-1}J_\theta F.
\end{equation}
当谱半径 $\rho(J_hF)<1$，逆矩阵有 Neumann 级数
\begin{equation}
 (I-J_hF)^{-1}=\sum_{k=0}^{\infty}(J_hF)^k,
\end{equation}
与无限展开的反向路径一致。动态有限深度模型通常不达到固定点，但此式说明稳定递归和可求导性的共同谱条件。

\section{自回归似然与 token 梯度}

对序列 $x_{1:T}$，因果语言模型按概率链式法则分解
\begin{equation}
 p_\theta(x_{1:T})=\prod_{t=1}^{T}p_\theta(x_t\mid x_{<t}).
\end{equation}
负对数似然
\begin{equation}
 \cL_{\rm NLL}=-\sum_{t=1}^{T}\log p_\theta(x_t\mid x_{<t}).
\end{equation}
若第 $t$ 步 logits 为 $z_t\in\R^{V}$，概率
$p_{tj}=e^{z_{tj}}/\sum_ke^{z_{tk}}$，真实 token 索引 $y_t$，则
\begin{equation}
 \ell_t=-z_{t,y_t}+\log\sum_ke^{z_{tk}},
\end{equation}
逐 logit 梯度
\begin{equation}
 \frac{\partial\ell_t}{\partial z_{tj}}
 =p_{tj}-\1\{j=y_t\}.
\end{equation}
Hessian 正是 softmax Jacobian
\begin{equation}
 \nabla_z^2\ell_t=\diag(p_t)-p_tp_t^T\succeq0.
\end{equation}
所以交叉熵关于 logits 凸，但关于深网参数并不凸。

对非 padding token 集合 $\mathcal I$，平均损失和困惑度
\begin{equation}
 \bar L=\frac1{|\mathcal I|}\sum_{t\in\mathcal I}\ell_t,
 \qquad \operatorname{PPL}=e^{\bar L}.
\end{equation}
因此 PPL 是真实 token 概率倒数的几何平均：
\begin{equation}
 \operatorname{PPL}=\left(
 \prod_{t\in\mathcal I}\frac1{p_\theta(x_t\mid x_{<t})}
 \right)^{1/|\mathcal I|}.
\end{equation}
不同 tokenizer 的 token 粒度不同，PPL 不能不经换算直接横向比较。

\section{旋转位置编码的相对位置恒等式}

对二维坐标对，定义旋转
\begin{equation}
 R(m\omega)=
 \begin{bmatrix}
 \cos(m\omega)&-\sin(m\omega)\\
 \sin(m\omega)&\cos(m\omega)
 \end{bmatrix}.
\end{equation}
位置 $m,n$ 的 query、key 分别为
$q_m=R(m\omega)q$、$k_n=R(n\omega)k$。内积
\begin{align}
 q_m^Tk_n
 &=q^TR(m\omega)^TR(n\omega)k\\
 &=q^TR((n-m)\omega)k.
\end{align}
第二步用到 $R(a)^T=R(-a)$ 和 $R(a)R(b)=R(a+b)$。
所以绝对位置旋转后的内积只依赖相对距离 $n-m$。

在 $d_h$ 维中对相邻坐标对使用频率
\begin{equation}
 \omega_j=\theta^{-2j/d_h},\qquad j=0,\ldots,d_h/2-1.
\end{equation}
低 $j$ 频率高、分辨近距离，高 $j$ 频率低、覆盖远距离。若推理长度远超训练长度，
相位 $m\omega_j$ 的外推和周期混叠会改变注意力几何；位置插值本质上是重新缩放这些相位。

\section{RMSNorm 的 Jacobian}

忽略可学习增益，RMSNorm
\begin{equation}
 y=\frac{x}{s},\qquad
 s=\sqrt{\frac1d x^Tx+\epsilon}.
\end{equation}
微分
\begin{equation}
 ds=\frac{x^Tdx}{d s},
\end{equation}
因此
\begin{align}
 dy&=\frac1s dx-\frac{x}{s^2}ds\\
 &=\left(\frac1sI-\frac{xx^T}{d s^3}\right)dx.
\end{align}
故
\begin{equation}
 J_{\rm RMS}=\frac1sI-\frac{xx^T}{d s^3}.
\end{equation}
与 LayerNorm 相比，它没有去均值投影
$I-d^{-1}\1\1^T$，因此不消除常数方向。加上逐坐标增益 $g$ 后，
Jacobain 左乘 $\diag(g)$。

谱上，对任意 $v\perp x$，$Jv=v/s$；沿 $x$ 方向
\begin{equation}
 Jx=\left(\frac1s-\frac{\|x\|^2}{d s^3}\right)x
 =\frac{\epsilon}{s^3}x.
\end{equation}
当 $\epsilon$ 很小，径向尺度方向被强烈压制，而切向方向按 $1/s$ 缩放。

\section{SwiGLU 前馈层的逐项梯度}

SwiGLU 可写
\begin{equation}
 a=xW_a,\qquad b=xW_b,\qquad
 h=\operatorname{swish}(a)\odot b,\qquad
 y=hW_o,
\end{equation}
其中
\begin{equation}
 \operatorname{swish}(a)=a\sigma(a).
\end{equation}
其导数
\begin{align}
 \frac d{da}[a\sigma(a)]
 &=\sigma(a)+a\sigma(a)(1-\sigma(a)).
\end{align}
令 $g_h=\partial\cL/\partial h$，则
\begin{align}
 g_a&=g_h\odot b\odot
 [\sigma(a)+a\sigma(a)(1-\sigma(a))],\\
 g_b&=g_h\odot\operatorname{swish}(a),\\
 \nabla_{W_a}\cL&=x^Tg_a,\qquad
 \nabla_{W_b}\cL=x^Tg_b,\\
 \nabla_x\cL&=g_aW_a^T+g_bW_b^T.
\end{align}
门分支 $b$ 同时调制激活和 $a$ 分支梯度，故两个投影不是可独立解释的普通 MLP。

\section{MoE 路由与负载}

对 token $x_i$，router logits $r_i=W_rx_i$，概率
\begin{equation}
 p_{ie}=\frac{e^{r_{ie}}}{\sum_{j=1}^{E}e^{r_{ij}}}.
\end{equation}
若 top-$k$ 专家集合 $S_i$，归一化门
\begin{equation}
 \widetilde p_{ie}=\frac{p_{ie}\1\{e\in S_i\}}
 {\sum_{j\in S_i}p_{ij}},
 \qquad
 y_i=\sum_{e\in S_i}\widetilde p_{ie}F_e(x_i).
\end{equation}
固定集合 $S_i$ 内可正常求导；集合边界的 arg-top-$k$ 不连续，几乎处处把未选专家梯度置零。

令软重要性
\begin{equation}
 P_e=\frac1N\sum_ip_{ie},
\end{equation}
硬负载
\begin{equation}
 F_e=\frac1N\sum_i\1\{e\in S_i\}/k.
\end{equation}
常见平衡项
\begin{equation}
 L_{\rm bal}=E\sum_eP_eF_e
\end{equation}
在均匀 $P_e=F_e=1/E$ 时等于 $1$。把硬 $F_e$ 停止梯度后，
router 梯度只通过 $P_e$，会降低当前过载专家的概率；它是训练启发式，不等同于严格容量约束。

若每专家容量为
\begin{equation}
 C_e=\left\lceil\frac{Nk}{E}\,c_{\rm cap}\right\rceil,
\end{equation}
超过容量的 token 必须丢弃或转送。即使平均负载均匀，局部批次方差仍可能溢出。

\section{低秩 KV 潜变量的吸收}

设 key/value 由共享潜变量
\begin{equation}
 c_t=W_cx_t\in\R^{d_c},\qquad
 k_t=W_kc_t,\qquad v_t=W_vc_t.
\end{equation}
注意力 logit
\begin{equation}
 q_s^Tk_t=q_s^TW_kc_t=(W_k^Tq_s)^Tc_t.
\end{equation}
因此解码时可把 $W_k$ 吸收到 query 投影，缓存 $c_t$ 而非完整 $k_t$。
若 value 聚合
\begin{equation}
 o_s=\sum_tp_{st}W_vc_t=W_v\left(\sum_tp_{st}c_t\right),
\end{equation}
$W_v$ 可放到加权和之后。缓存从每 token 的 $d_k+d_v$ 降为 $d_c$，
前提是位置编码等操作不破坏这些线性交换关系；对只作用在部分 key 维度的 RoPE，需把旋转部分另行缓存或重构。

\section{多 token 预测的联合似然}

若同一隐藏状态 $h_t$ 用 $K$ 个头预测 $x_{t+1},\ldots,x_{t+K}$，
\begin{equation}
 \cL_t=\sum_{k=1}^{K}\lambda_k
 [-\log p_{\theta,k}(x_{t+k}\mid h_t)].
\end{equation}
共享骨干梯度
\begin{equation}
 \nabla_{h_t}\cL_t=\sum_{k=1}^{K}\lambda_k g_{t,k}.
\end{equation}
其平方范数精确展开
\begin{equation}
 \left\|\sum_k\lambda_kg_{t,k}\right\|^2
 =\sum_k\lambda_k^2\|g_{t,k}\|^2
 +2\sum_{j<k}\lambda_j\lambda_k g_{t,j}^Tg_{t,k}.
\end{equation}
交叉内积为负时存在任务冲突；辅助头退火相当于训练后期逐渐去掉这些交叉项，
使最终目标回到下一 token 似然。

\end{document}
